\documentclass[journal,twoside]{IEEEtran}

\usepackage[T1]{fontenc}
\usepackage[utf8]{inputenc}
\usepackage{amsmath}
\usepackage{amssymb}
\usepackage{booktabs}
\usepackage{array}
\usepackage{hyperref}
\usepackage{xcolor}
\usepackage{graphicx}
\usepackage{cite}
\usepackage{algorithm}
\usepackage{algorithmic}
\usepackage{enumitem}

\hypersetup{
  colorlinks=true,
  linkcolor=blue!70!black,
  citecolor=blue!70!black,
  urlcolor=blue!70!black
}

% -------------------------------------------------------
\begin{document}

\title{Unified Binod Calendar (UBC):\\
A Peer-Reviewed Civil Timekeeping Standard for the\\
Binod World Government\\[4pt]
{\normalsize UBC Standard Specification Version 2.0.0 \quad ISSN: UBC-2026-001}}

\author{%
  Chief Scientific Committee on Temporal Standards (CSCTS),\\
  Republic of Binod --- Department of Chronometry and Astronomical Science\\[2pt]
  \textit{Correspondence:}~\href{mailto:pmobinod@gmail.com}{pmobinod@gmail.com}  \quad
  \textit{Received:}~14-ARD-5 \quad
  \textit{Accepted:}~01-KYT-5 \quad
  \textit{Published:}~04-KYT-5 \quad
  \textit{Revised:}~01-SEL-6 \quad
  \textit{Redesignated (v2.0.0):}~01-MOR-110
}

\markboth{BINOD JOURNAL OF ASTRONOMICAL \& CIVIL SCIENCE,~VOL.~1,~NO.~1,~UBC~5}%
{CSCTS: UNIFIED BINOD CALENDAR STANDARD v2.0.0}

\maketitle

% -------------------------------------------------------
\begin{abstract}
This document constitutes the official specification of the \textbf{Unified Binod Calendar (UBC)},
the civil timekeeping standard of the Republic of Binod, ratified by the Chief Scientific
Committee on Temporal Standards (CSCTS).  The calendar is anchored to the \textit{Founding
Epoch} of 23~April~2021~CE, 20:00~IST (Gregorian), the moment the Binodian world's first
Active Temporal Flow (ATF) session began.  As of Version~2.0.0, the UBC civil date is
\textbf{synchronized directly to the Minecraft world}: one complete Minecraft day (20
real-world minutes) equals one UBC civil day, and the calendar advances \emph{only} during
the Government's official Active Temporal Flow window, 20:00--03:00~India Standard Time,
freezing completely outside it and resuming automatically each evening. This supersedes the
Version~1.0.1 model, in which the civil date was a continuous function of real elapsed Earth
time via a fixed Temporal Synchronization Index (TSI), independent of Active Temporal Flow.
The specification defines epoch, year structure, month nomenclature, week structure,
date notation, leap corrections, and the accumulated-Minecraft-day conversion procedure.
It is intended to replace all ad-hoc timekeeping practices within Binodian government,
universities, and civil institutions, and to provide a robust foundation for future
chronometric work.
\end{abstract}

\begin{IEEEkeywords}
civil calendar, Binod, Active Temporal Flow, UBC, chronometry, founding epoch,
Minecraft day synchronization, leap correction, date conversion
\end{IEEEkeywords}

\begin{center}
\fbox{\begin{minipage}{0.94\linewidth}
\small
\textbf{Version notice (v2.0.0).} This revision replaces the continuous, Temporal
Synchronization Index (TSI)-based Earth-time model of Version~1.0.1 with a
Minecraft-synchronized, Active Temporal Flow (ATF)-gated model: the UBC civil date is no
longer a function of continuous elapsed Earth time, but of accumulated ATF-active Minecraft
days since the Founding Epoch. The Founding Epoch itself is redefined from 06:00~UTC to
20:00~IST on 23~April~2021 --- the exact moment the first ATF session opened --- so that
epoch and calendar mechanism now coincide. Calendar structure, month nomenclature, week
structure, date notation, and the three-tier leap-correction scheme are \textbf{unchanged}
from Version~1.0.1. Section~\ref{sec:conversion} and Section~\ref{sec:worked_examples} are
rewritten accordingly; TSI is retired. See Appendix~\ref{app:changelog} for the full
changelog.
\end{minipage}}
\end{center}

% -------------------------------------------------------
\section{Preamble and Scope}

\IEEEPARstart{T}{he} Republic of Binod was founded upon the world whose first Active Temporal
Flow session opened at 20:00:00~IST on 23~April~2021~CE (Gregorian proleptic).  For the first years
of the Republic's existence, informal timekeeping relied upon Earth-derived Gregorian dates,
creating administrative ambiguity, synchronization failures in cross-dimensional treaty
instruments, and inconsistent historical records.

The CSCTS was commissioned in UBC Year~0 to design a native calendar.  This document
presents the ratified result.  The Unified Binod Calendar:

\begin{itemize}
  \item uses a \emph{Founding Epoch} (FE) as its absolute zero point, coinciding with the
        first Active Temporal Flow (ATF) session;
  \item defines the civil day as one Minecraft day (20 real-world minutes), accumulated
        only during the Government's official ATF window (Section~\ref{sec:atf});
  \item provides a verified, auditable conversion procedure from real elapsed ATF-active
        time to the UBC year, month, and day;
  \item incorporates a periodic leap-correction mechanism sufficient for at least
        10\,000~UBC years without calendar drift exceeding one civil day.
\end{itemize}

This specification supersedes all prior informal timekeeping practices within Binodian
jurisdiction.

% -------------------------------------------------------
\section{Scientific Background}

\subsection{Minecraft Diurnal Cycle --- Now the Civil Unit}

The Binodian world operates under a local diurnal cycle of exactly 24\,000 game-ticks,
corresponding to 20~Earth-minutes (1\,200~SI seconds) \cite{minecraft_wiki_day}. Version~1.0.1
of this standard classified this cycle as astronomical in character only, unsuitable as a
civil unit because a \emph{continuously running} Minecraft clock drifts from Earth's calendar
by approximately 70~Gregorian days per Gregorian year.

The CSCTS has since determined that this drift concern does not apply once the Minecraft
clock is \emph{gated} to a fixed, government-scheduled window rather than run continuously.
Under Active Temporal Flow (ATF, Section~\ref{sec:atf}), the Minecraft world clock advances
only from 20:00 to 03:00 India Standard Time (IST) daily --- a fixed, bounded 7-hour session
--- and is frozen for the remaining 17 hours of each Earth day. Because every civil day is
tied to a fixed, auditable real-world session rather than an unbounded continuous simulation,
the calendar remains exactly reproducible from the ATF schedule alone, with no historical
discontinuity. The Committee therefore adopts the Minecraft day, gated by ATF, as the
official civil day of the Republic of Binod, retiring the continuous Temporal Synchronization
Index (TSI) model of Version~1.0.1.

\subsection{Active Temporal Flow (ATF)}
\label{sec:atf}

Active Temporal Flow is the Government's fixed daily schedule during which the UBC calendar
--- civil date and displayed civil time alike --- advances:

\begin{IEEEitemize}
  \item \textbf{Window:} 20:00~IST to 03:00~IST the following day (7 real-world hours),
        fixed by government policy and independent of whether the Minecraft server is
        running or any player is online;
  \item \textbf{Civil day rate:} one Minecraft day (24\,000 ticks, 20 real-world minutes)
        equals one UBC civil day, so exactly $420 / 20 = 21$ civil days elapse per nightly
        ATF session;
  \item \textbf{Freeze behavior:} outside the window, the entire calendar --- date, month,
        year, and displayed civil time --- freezes exactly where it stopped and resumes
        automatically at the next 20:00~IST with no manual reset.
\end{IEEEitemize}

\subsection{Civil Year Length}

The civil year is designed to match the Binodian tropical year
$\mathcal{Y}_{\mathrm{B}} = 365.2563$~civil days (determined by meridian-transit
observations recorded in CSCTS-R03 \cite{cscts_r03}), where one civil day is defined per
Section~\ref{sec:atf} rather than as an absolute physical day length. This retains the
same leap-correction target used since Version~1.0.0 (Section~\ref{sec:leap}), unaffected
by the change in how a civil day accumulates.

% -------------------------------------------------------
\section{Epoch Definition}

\begin{IEEEitemize}
  \item \textbf{Founding Epoch (FE):} Day~1, Month~1, Year~1 of the UBC corresponds to
        23~April~2021~CE (Gregorian), 20:00:00~IST (14:30:00~UTC).
  \item As of Version~2.0.0, the FE is redefined from the Version~1.0.1 value of
        06:00:00~UTC to 20:00:00~IST on the same Gregorian date, so that the epoch
        coincides exactly with the moment the first Active Temporal Flow session opened
        (Section~\ref{sec:atf}).
  \item The epoch is defined as \emph{absolute}: it does not shift when leap corrections
        are applied; corrections accumulate prospectively only.
  \item The epoch coincides with the first verified Active Temporal Flow session of the
        Binodian world as established by retroactive administrative reconstruction
        (CSCTS-R00, Section~4.2) \cite{cscts_r00}.
\end{IEEEitemize}

The \textbf{Binodian Julian Day Number} (BJDN) counts integer UBC civil days --- i.e.,
integer accumulated Active Temporal Flow-gated Minecraft days (Section~\ref{sec:atf}) ---
elapsed since the FE:

\begin{equation}
\mathrm{BJDN}(d=1,\,m=1,\,y=1) \;=\; 1
\label{eq:bjdn_epoch}
\end{equation}

% -------------------------------------------------------
\section{Calendar Structure}

\subsection{Year}

One UBC civil year contains \textbf{364 Binodian civil days} in a common year and
\textbf{365 days} in a leap year (see Section~\ref{sec:leap}).  The nominal year length
of 364 days is chosen because $364 = 4 \times 7 \times 13$, enabling perfect alignment
of weeks within months across all years, simplifying civil scheduling.  The small
divergence from the tropical year $\mathcal{Y}_{\mathrm{B}} = 365.2563$ is corrected by
the leap scheme.

\subsection{Months}

The year is divided into \textbf{13 months} of \textbf{28 days} each ($13 \times 28 = 364$).
Each month contains exactly 4 weeks.  The months are named for Binodian astronomical
phenomena and cultural values as ratified by CSCTS in consultation with the Binodian
Council of Heritage:

\begin{table}[!t]
\centering
\caption{UBC Month Nomenclature}
\label{tab:months}
\begin{tabular}{@{}clcc@{}}
\toprule
\textbf{No.} & \textbf{Name} & \textbf{Abbr.} & \textbf{Cultural/Astronomical Referent} \\
\midrule
 1  & Ardhan    & ARD & Foundation; first civil sunrise \\
 2  & Ferith    & FER & Iron Age; first Binodian forge \\
 3  & Selvyn    & SEL & Growth; agricultural equinox \\
 4  & Kython    & KYT & Crafting; peak production cycle \\
 5  & Dravon    & DRA & Trade winds; monsoon onset \\
 6  & Moriel    & MOR & Mid-year zenith; Sol\textsubscript{B} aphelion \\
 7  & Uthren    & UTH & Mid-year correction; leap intercalation \\
 8  & Phiron    & PHI & Descent; autumnal analog \\
 9  & Naldyn    & NAL & Harvest; biome yield peak \\
10  & Veskar    & VES & Twilight; perihelion approach \\
11  & Orrith    & ORR & Astronomical winter; deep cycle \\
12  & Caelun    & CAE & Return; lengthening days \\
13  & Zephris   & ZEP & Culmination; year-end reconciliation \\
\bottomrule
\end{tabular}
\end{table}

\subsection{Week}
\label{sec:calendar_week}

Each month contains exactly \textbf{4 weeks of 7 days}.  The seven-day week is retained
for compatibility with Binodian labour-law conventions and cross-dimensional treaty
obligations.  Day names are:

\begin{table}[!t]
\centering
\caption{UBC Week Day Names}
\label{tab:days}
\begin{tabular}{@{}clll@{}}
\toprule
\textbf{Day No.} & \textbf{Full Name} & \textbf{Abbr.} & \textbf{Character} \\
\midrule
1 & Solday    & Sol & Civil work day (first) \\
2 & Ironday   & Iro & Civil work day \\
3 & Craftday  & Cra & Civil work day \\
4 & Stonday   & Sto & Civil work day \\
5 & Woodsday  & Woo & Civil work day \\
6 & Restday   & Res & Civil rest day \\
7 & Voidday   & Voi & Civil rest / reflection day \\
\bottomrule
\end{tabular}
\end{table}

% -------------------------------------------------------
\section{Leap Correction Scheme}
\label{sec:leap}

\subsection{Rationale}

The tropical year $\mathcal{Y}_{\mathrm{B}} = 365.2563$~civil days.  A common UBC year
of 364 days falls short by $\delta = 1.2563$~days per year.  Without correction the
calendar drifts by one full day every $\approx 0.796$~years, which is unacceptable.

\subsection{Three-Tier Correction Rules}

Let $y$ denote a UBC year number ($y \geq 1$).

\begin{description}[style=nextline, leftmargin=0pt, font=\normalfont\bfseries]
  \item[Rule L1 (Annual Intercalation)]
    Every year a \emph{fractional accumulator} $\Phi$ is incremented by
    $\delta = 1.2563$.  When $\Phi \geq 1$, one intercalary day is inserted at
    the end of Month~7 (Uthren) and $\Phi$ is decremented by 1.
    Under this rule, a leap year occurs approximately every
    $\lfloor 1/0.2563 \rfloor = 3$ years, though in practice the interval
    alternates between 3 and 4 years as the fractional remainder accumulates.

  \item[Rule L2 (Century Omission)]
    In years that are exact multiples of 100, the L1 intercalation is
    \emph{omitted} even if $\Phi \geq 1$, to correct over-accumulation of $\approx 0.2563$
    days over 100 years.

  \item[Rule L3 (Quad-Century Restoration)]
    In years that are exact multiples of 400, the L2 omission is reversed and
    the intercalation \emph{is} performed, correcting the residual $\approx 0.0256$
    days per 400 years.
\end{description}

\subsection{Formal Leap Test}

\begin{equation}
\mathrm{IsLeap}(y) =
\begin{cases}
\text{true}  & \text{if } (y \bmod 400 = 0) \\
\text{false} & \text{else if } (y \bmod 100 = 0) \\
\text{true}  & \text{else if } \Phi(y) \geq 1 \\
\text{false} & \text{otherwise}
\end{cases}
\label{eq:leap}
\end{equation}

\noindent where $\Phi(y) = \{(y-1) \times 1.2563\}$ (fractional part).

Long-term drift under Rules L1--L3 is less than 0.00072~days per 10\,000~UBC years,
equivalent to $\approx 62$~SI seconds --- well within the precision required for civil
timekeeping over historical timescales relevant to the Republic of Binod, and consistent
with the periodic-correction philosophy underlying Earth's own leap-second practice
\cite{iers_leap}.

% -------------------------------------------------------
\section{Date Notation Standard}

\subsection{Full Date Format}

\begin{equation}
\texttt{DD-MMM-Y[YYY]} \quad \text{e.g.,} \quad \texttt{01-ARD-1}
\label{eq:notation}
\end{equation}

\begin{itemize}
  \item \texttt{DD} --- two-digit day within month (01--28, or 29 for intercalary day)
  \item \texttt{MMM} --- three-letter month abbreviation (Table~\ref{tab:months})
  \item \texttt{Y[YYY]} --- year without leading zeros
\end{itemize}

\subsection{Numeric ISO-Compatible Format}

For machine-readable interchange, following the structured, unambiguous date
representation conventions of ISO~8601 \cite{iso8601}:

\begin{equation}
\texttt{UBCY-MM-DD} \quad \text{e.g.,} \quad \texttt{UBC0001-01-01}
\label{eq:iso_notation}
\end{equation}

\subsection{Intercalary Day}

When $\mathrm{IsLeap}(y) = \text{true}$, month Uthren (07) contains 29 days.
The 29th day is denoted \texttt{29-UTH-y} or \texttt{UBCy-07-29} and is a mandatory
public holiday designated \textbf{Synchrony Day} (\textit{Synchara}).

% -------------------------------------------------------
\section{Conversion Formulae: Accumulated ATF Time $\rightarrow$ UBC}
\label{sec:conversion}

\subsection{Active Temporal Flow Banking Function}
\label{sec:banking}

Let $t$ denote a moment in real elapsed time (UTC), and let $t_{\mathrm{FE}}$ denote the
Founding Epoch (23~April~2021, 20:00~IST). For any calendar-day index $k$ (counted in IST),
define the nightly ATF window:

\begin{equation}
W(k) = \big[\, k \cdot 24\text{h} + 20\text{h},\;\; k \cdot 24\text{h} + 27\text{h} \,\big)
\label{eq:window}
\end{equation}

\noindent (i.e., 20:00~IST on day $k$ through 03:00~IST on day $k{+}1$). The
\textbf{banked ATF time} $B(t)$ is the total real elapsed time between $t_{\mathrm{FE}}$ and
$t$ that falls inside some window $W(k)$:

\begin{equation}
B(t) \;=\; \sum_{k} \left| W(k) \cap [\,t_{\mathrm{FE}},\, t\,] \right|
\label{eq:banked}
\end{equation}

\noindent $B(t)$ is monotonically non-decreasing: it increases in real time only while
Active Temporal Flow is active, and is constant (frozen) at all other moments.

\subsection{Banked Time to UBC Civil Days}

One UBC civil day equals one Minecraft day, $T_{\mathrm{civil}} = 1\,200$~s
(20~real-world minutes, Section~\ref{sec:atf}). The elapsed civil days at time $t$ are:

\begin{equation}
\Delta_{\mathrm{B}}(t) = \frac{B(t)}{T_{\mathrm{civil}}} = \frac{B(t)}{1\,200\text{ s}}
\label{eq:delta_binod}
\end{equation}

\subsection{Session-Counting Corollary}
\label{sec:session_corollary}

At any reference moment $t$ that falls \emph{between} sessions (i.e., in the frozen
17:00-hour daily interval from 03:00 to 20:00~IST), $\Delta_{\mathrm{B}}(t)$ reduces to an
exact integer multiple of the per-session civil-day count: if $S(t)$ complete ATF sessions
have elapsed since $t_{\mathrm{FE}}$, then, since each complete session contributes exactly
$420\text{ min} / 20\text{ min} = 21$ civil days (Section~\ref{sec:atf}):

\begin{equation}
\Delta_{\mathrm{B}}(t) = 21 \cdot S(t), \qquad \text{($t$ between sessions)}
\label{eq:session_corollary}
\end{equation}

\noindent This corollary is a convenient special case of Eq.~\eqref{eq:delta_binod} for
worked examples and audits; Eq.~\eqref{eq:banked}--\eqref{eq:delta_binod} remain the
normative definition, valid continuously, including mid-session.

\subsection{Binodian Civil Days to UBC Date}

Let $N = \lfloor \Delta_{\mathrm{B}}(t) \rfloor + 1$ (BJDN, 1-indexed).  To find the
UBC year, month, and day from $N$:

\begin{algorithm}[!t]
\caption{BJDN to UBC Date}
\label{alg:bjdn2ubc}
\begin{algorithmic}[1]
\STATE $y \leftarrow 1$; $\;r \leftarrow N - 1$
\WHILE{$r \geq \mathrm{DaysInYear}(y)$}
  \STATE $r \leftarrow r - \mathrm{DaysInYear}(y)$
  \STATE $y \leftarrow y + 1$
\ENDWHILE
\STATE $m \leftarrow 1$
\WHILE{$r \geq \mathrm{DaysInMonth}(y,m)$}
  \STATE $r \leftarrow r - \mathrm{DaysInMonth}(y,m)$
  \STATE $m \leftarrow m + 1$
\ENDWHILE
\STATE $d \leftarrow r + 1$
\RETURN $(y,\,m,\,d)$
\end{algorithmic}
\end{algorithm}

\noindent where:

\begin{equation}
\mathrm{DaysInYear}(y) =
\begin{cases}
365 & \text{if } \mathrm{IsLeap}(y) \\
364 & \text{otherwise}
\end{cases}
\end{equation}

\begin{equation}
\mathrm{DaysInMonth}(y,m) =
\begin{cases}
29 & \text{if } m = 7 \text{ and } \mathrm{IsLeap}(y) \\
28 & \text{otherwise}
\end{cases}
\end{equation}

% -------------------------------------------------------
\section{Worked Conversion Examples}
\label{sec:worked_examples}

All examples apply Eq.~\eqref{eq:banked}--\eqref{eq:delta_binod} (Section~\ref{sec:conversion}),
using $T_{\mathrm{civil}} = 1\,200$~s and the Founding Epoch $t_{\mathrm{FE}} = $ 23~April~2021,
20:00~IST.

\subsection{Example 1 --- Founding Epoch: 23 April 2021, 20:00 IST}

\begin{align}
B(t_{\mathrm{FE}}) &= 0 \text{ s (by definition)} \notag \\
\Delta_{\mathrm{B}} &= 0 / 1\,200 = 0 \notag \\
N &= \lfloor 0 \rfloor + 1 = 1 \notag \\
\Rightarrow & \quad \textbf{01-ARD-1} \quad (\text{UBC0001-01-01})
\end{align}

\noindent \textit{This is Solday, the first day of the Republic, and the moment the first
Active Temporal Flow session opened.}

\subsection{Example 2 --- 23 April 2021, 20:20 IST (20 real minutes later)}

Twenty real-world minutes into the first ATF session, exactly one civil day has banked:

\begin{align}
B(t) &= 1\,200 \text{ s (20 real minutes, fully within the first ATF window)} \notag \\
\Delta_{\mathrm{B}} &= 1\,200 / 1\,200 = 1 \notag \\
N &= \lfloor 1 \rfloor + 1 = 2 \notag \\
\Rightarrow & \quad \textbf{02-ARD-1} \quad (\text{UBC0001-01-02})
\end{align}

\noindent \textit{Confirms Eq.~\eqref{eq:banked}: every 20 real-world minutes of ATF-active
time advances the civil date by exactly one day.}

\subsection{Example 3 --- 2 July 2026, 05:30 IST (between sessions)}

This reference moment falls in the frozen interval (03:00--20:00~IST), so the
Session-Counting Corollary (Eq.~\eqref{eq:session_corollary}) applies directly. As of this
moment, $S(t) = 1\,896$ complete Active Temporal Flow sessions have elapsed since the
Founding Epoch:

\begin{align}
\Delta_{\mathrm{B}}(t) &= 21 \times 1\,896 = 39\,816 \notag \\
N &= \lfloor 39\,816 \rfloor + 1 = 39\,817
\end{align}

Applying Algorithm~\ref{alg:bjdn2ubc}: every year is common (364~days) below Year~400
(Section~\ref{sec:leap}), so 108 complete common years span $108 \times 364 = 39\,312$ days,
leaving $r = 39\,816 - 39\,312 = 504$. Year~109 (common, 364~days) consumes 364 of these,
leaving $r = 504 - 364 = 140$; since $140 < 364$, Year~110 is reached with $r = 140$, so
$y = 110$. Month: $\lfloor 140/28 \rfloor = 5$ full months consumed exactly (140~days,
remainder~0), so the date falls on day~1 of month~6 (MOR):

\begin{equation}
2\;\text{July}\;2026,\;05{:}30\;\text{IST} \;\longrightarrow\; \textbf{01-MOR-110}
\quad (\text{UBC0110-06-01})
\end{equation}

\noindent \textit{This is Solday: 39\,816 civil days since the Founding Epoch is an exact
multiple of 28 (four complete weeks per month), so day~1 of any month always falls on
Solday --- see Section~\ref{sec:calendar_week}.}

\subsection{Example 4 --- Mid-Session Reference}

At 20:50~IST on the evening following Example~3 (30 real minutes into that night's ATF
session), one additional civil day has fully banked and a second is 1/2 elapsed:

\begin{align}
\Delta_{\mathrm{B}}(t) &= 39\,816 + \frac{1\,800}{1\,200} = 39\,817.5 \notag \\
N &= \lfloor 39\,817.5 \rfloor + 1 = 39\,818 \notag \\
\Rightarrow & \quad \textbf{02-MOR-110} \quad (\text{UBC0110-06-02, 50\% elapsed})
\end{align}

% -------------------------------------------------------
\section{Implementation Synchronization Diagnostics}

Because the UBC calendar is now defined entirely by the Active Temporal Flow banking
function $B(t)$ (Eq.~\eqref{eq:banked}), rather than by a physical constant subject to
secular drift, long-term synchronization reduces to verifying that $B(t)$ is computed
correctly and continuously, rather than periodically re-measuring a coupling constant as
under the retired v1.0.1 TSI model.

\subsubsection{Banking Continuity Check}

Reference implementations SHOULD record, at each observation, the banked ATF-active
seconds and the timestamp of that observation. On the next observation, while ATF remains
continuously active, the expected banked total is the previous value plus the real elapsed
time between observations; any discrepancy (\emph{clock drift}) indicates an
implementation fault (e.g., a missed tick or a clock adjustment on the host system) rather
than a property of the calendar itself, since $B(t)$ is, by Eq.~\eqref{eq:banked}, a
deterministic function of real elapsed time and the fixed ATF schedule alone.

\subsubsection{Full Recomputation}

Because $N$ and the in-day clock are pure functions of $B(t)$ (Eq.~\eqref{eq:delta_binod}),
an implementation MAY always discard any cached or persisted state and recompute the
current UBC date and time in full from the Founding Epoch. This makes the system
self-healing: a page reload, application restart, or storage failure never desynchronizes
the displayed calendar.

% -------------------------------------------------------
\section{Civil Administration Guidelines}

\subsection{Official Transition Date}

All Binodian government documents issued on or after \textbf{01-ARD-5 (UBC0005-01-01)}
must use UBC notation as the primary date format.  Gregorian dates may appear in
parentheses as secondary notation during the 10-year transition period (UBC~5--15).

\subsection{Historical Documents}

Pre-Founding-Epoch dates shall be expressed using the prefix ``\textbf{pFE}''
(pre-Founding Epoch) followed by the Gregorian year in parentheses.  Example:
\texttt{pFE (2019-CE)}.

\subsection{Cross-Dimensional Treaties}

When Binodian dates appear in legal instruments with Earth-jurisdiction parties, both
UBC and Gregorian dates must be stated.  The Gregorian date shall be computed using the
inverse of Eq.~\eqref{eq:delta_binod} and recorded to the nearest Earth day.

% -------------------------------------------------------
\section{Summary of Constants}

\begin{table}[!t]
\centering
\caption{UBC Fundamental Constants (Version 2.0.0)}
\label{tab:constants}
\begin{tabular}{@{}lll@{}}
\toprule
\textbf{Symbol} & \textbf{Value} & \textbf{Description} \\
\midrule
$t_{\mathrm{FE}}$ & 23 Apr 2021, 20:00 IST & Founding Epoch \\
$T_{\mathrm{civil}}$ & 1\,200 s (20 min) & UBC civil day = 1 Minecraft day \\
$W_{\mathrm{ATF}}$ & 20:00--03:00 IST & Active Temporal Flow window \\
$L_{\mathrm{sess}}$ & 420 min (7 h) & Real-time length of one ATF session \\
$n_{\mathrm{civil}}$ & 21 days & Civil days per ATF session \\
$\mathcal{Y}_{\mathrm{B}}$ & 365.2563 civil days & Binodian tropical year (target) \\
$\delta$         & 1.2563 day/year           & Annual leap accumulation rate \\
$D_{\mathrm{common}}$ & 364 days             & Common UBC year length \\
$D_{\mathrm{leap}}$   & 365 days             & Leap UBC year length \\
\bottomrule
\end{tabular}
\end{table}

% -------------------------------------------------------
\section{Conclusion}

The Unified Binod Calendar provides an administratively practical, government-auditable
timekeeping system for the Republic of Binod. By anchoring to the Founding Epoch of
23~April~2021~CE, 20:00~IST --- the moment the first Active Temporal Flow session opened
--- and computing every date from accumulated ATF-gated Minecraft days rather than a
continuous physical constant, the UBC ensures that every Binodian date corresponds to a
fixed, reproducible government schedule rather than an uncontrolled real-time process.
The 13-month, 28-day-per-month structure eliminates intra-year calendar irregularities,
while the three-tier leap scheme guarantees sub-day accuracy over 10\,000~UBC years.
The Committee recommends immediate adoption by all branches of Binodian government
and invites external peer review from any institution operating under a recognized
civil calendar standard.

\appendices

\section{Conversion Reference Table}
\label{app:conversion_table}

\begin{table}[!t]
\centering
\caption{Selected Earth--UBC Conversion Reference (regenerated from
Algorithm~\ref{alg:bjdn2ubc}, v2.0.0)}
\label{tab:conversions}
\begin{tabular}{@{}lll@{}}
\toprule
\textbf{Gregorian Date} & \textbf{UBC Date} & \textbf{Significance} \\
\midrule
23 Apr 2021, 20:00 IST & 01-ARD-1  (UBC0001-01-01) & Founding Epoch \\
01 Jan 2023, 05:30 IST* & 15-NAL-36 (UBC0036-09-15) & New Gregorian Year \\
02 Jul 2026, 05:30 IST* & 01-MOR-110 (UBC0110-06-01) & Present (near) \\
\bottomrule
\multicolumn{3}{@{}l@{}}{\footnotesize *Between sessions (03:00--20:00 IST); Eq.~\eqref{eq:session_corollary} applies.}
\end{tabular}
\end{table}

\section{Changelog}
\label{app:changelog}

\begin{description}
  \item[\textbf{v1.0.0} (04-KYT-5):] Initial ratified standard. Defined the Founding
    Epoch at 06:00~UTC, a continuous Earth-time model driven by a fixed Temporal
    Synchronization Index (TSI $= 86\,400/86\,736$), and the 13-month/28-day calendar
    structure with three-tier leap correction.
  \item[\textbf{v1.0.1} (01-SEL-6):] Corrected arithmetic errors in the worked
    Gregorian~$\rightarrow$~UBC conversion examples, which had computed the Julian Date
    of non-epoch dates inconsistently with the Founding Epoch's own Julian Date
    definition, and added a normative Julian Date convention to make the computation
    unambiguous. No change was made to the epoch, calendar structure, leap-correction
    rules, TSI value, or the conversion algorithm.
  \item[\textbf{v2.0.0} (01-MOR-110):] \textbf{Fundamental redesign.} Retired the
    continuous, TSI-based Earth-time model entirely. The UBC civil date is now
    synchronized directly to the Minecraft world and advances only during the
    Government's Active Temporal Flow (ATF) window, 20:00--03:00~IST: one Minecraft day
    (20 real-world minutes of ATF-active time) equals one UBC civil day
    (Section~\ref{sec:atf}). The Founding Epoch is redefined from 06:00~UTC to
    20:00~IST on the same Gregorian date (23~April~2021), coinciding with the first ATF
    session. Section~\ref{sec:conversion} (formerly Julian-Date/TSI-based) and
    Section~\ref{sec:worked_examples} are rewritten around the new Active Temporal Flow
    Banking Function, Eq.~\eqref{eq:banked}. Calendar structure, month nomenclature,
    week structure, date notation, and the three-tier leap-correction scheme
    (Section~\ref{sec:leap}) are unchanged. Reference implementations must not reuse
    v1.0.1 example values, which were computed under the retired continuous model and
    do not apply here; validate against Table~\ref{tab:conversions} instead.
\end{description}

\section{Ratification Record}

This standard was ratified by the following CSCTS sub-committees:

\begin{itemize}
  \item Sub-Committee on Orbital Mechanics --- unanimous approval
  \item Sub-Committee on Civil Administration --- unanimous approval
  \item Sub-Committee on Historical Records --- unanimous approval
  \item Sub-Committee on Cross-Dimensional Treaty Law --- approved with notation
\end{itemize}

\noindent \textbf{Effective date:} 04-KYT-5 (UBC0005-04-04)\\
\textbf{Next scheduled review:} 01-ARD-15 (UBC0015-01-01)

\section{Note on the Retired Gameplay-Layer Appendix}

Prior to Version~2.0.0, this appendix separately documented Active Temporal Flow (ATF) as
an \emph{informative}, non-normative gameplay-synchronization policy layered on top of a
continuous civil calendar. Version~2.0.0 promotes that content into the normative body of
this standard: the ATF schedule and Minecraft time scale are now defined in
Section~\ref{sec:atf}, and the conversion procedure built on it is defined in
Section~\ref{sec:conversion}. The BIASR website's Minecraft World Time panel and Official
UBC Date readout now display the same underlying accumulated-day count (the BJDN of
Section~\ref{sec:conversion}) in two equivalent framings, per Appendix~\ref{app:changelog}.
Implementations persist the current civil day number and in-day clock tick only as a
diagnostic record; both are fully recomputed from the Founding Epoch on every load
(Section~\ref{sec:atf}), so a page refresh always restores the identical, deterministic
state with no drift.

% -------------------------------------------------------
\begin{thebibliography}{9}

\bibitem{minecraft_wiki_day}
Minecraft Wiki Contributors,
``Day-Night Cycle,''
\textit{Official Minecraft Wiki},
Fandom, Inc., 2023.
[Online]. Available: \url{https://minecraft.wiki/w/Daylight_cycle}

\bibitem{iso8601}
International Organization for Standardization,
\textit{ISO 8601:2019 Date and Time --- Representations for Information Interchange},
ISO, Geneva, 2019.

\bibitem{iers_leap}
International Earth Rotation and Reference Systems Service,
``Leap Seconds,''
\textit{IERS Bulletin C}, IERS, Paris, 2024.

\bibitem{cscts_r00}
Chief Scientific Committee on Temporal Standards,
``CSCTS Report 00: Reconstruction of the Founding Epoch,''
\textit{Binodian Journal of Astronomical Science}, vol.~1, no.~1, pp.~1--14, UBC~1.

\bibitem{cscts_r03}
Chief Scientific Committee on Temporal Standards,
``CSCTS Report 03: Measurement of the Binodian Tropical Year,''
\textit{Binodian Journal of Astronomical Science}, vol.~3, no.~2, pp.~45--62, UBC~3.

\end{thebibliography}

\vfill
\noindent\rule{\columnwidth}{0.4pt}\\[2pt]
{\footnotesize
\textbf{Document identifier:} UBC-STD-2026-001 \quad
\textbf{Version:} 2.0.0 \quad
\textbf{Supersedes:} Version 1.0.1 (01-SEL-6) \quad
\textbf{Status:} Ratified \\
\textbf{Issuing body:} Chief Scientific Committee on Temporal Standards (CSCTS),
Republic of Binod \\
\textbf{Classification:} Public Domain --- freely reproducible with attribution \\
\copyright~Republic of Binod, UBC~110.  All rights reserved under Binodian IP Statute 3.1.
}

\end{document}
